\section{基于二维线程块与线程的矩阵加法}

欢迎各位来到本课程中的重要文章，我们将开始第二阶段的讲解。在所有先前的应用中，我们专注于使用一维线程和线程块。我们将此应
用在一个简单的例子中，该例子是将两个向量相加，并将结果存储在第三个向量中，并且在每个网格中使用一维线程块。

今天我们需要更深入一层，或者说进入本课程的下一阶段，即我们在每个线程块中使用二维线程。由于我们需要在一个实际例子中应用这
一点，我在此添加了这个标题：二维矩阵加法。今天的主要目的是将两个矩阵相加，并将结果存储在第三个矩阵中。也就是说不再使用向
量，明白吗？让我们在这里复习一下。

在之前的讨论中我告诉过你们，我们在所有 CUDA 应用中使用线程块大小和网格大小，对吧？线程块大小在 CUDA代码中被称
为 \texttt{blockDim.x}。这一点很重要，因为我们需要在 CUDA应用的方程式中使用这个标识符。并且它也
可以标识为每个线程块的线程数。

记住，线程块大小等于 X 维度的 \texttt{blockDim}。因为我们之前只使用了一个维度，对于 Ampere架
构，我们在这种情况下可以使用的最大值是每个线程块 1024 个线程。这也就等于每个线程块的线程数。对于 Ampere架构
我们可以使用的最大值是 1024。然后我们还有网格大小，网格大小被称为\texttt{gridDim.x}。我们需要这个
是因为有时我们需要在 CUDA代码中使用它，这里最大值非常大，但我这里只是举个例子，所以这不是最大值。例如我们可以使用每
个网格 1024个线程块，那样的话我们就可以计算整个应用的线程总数，而不是每个线程块的线程数。你只需将 1024 乘以
1024就能得到整个应用的总线程数。这是假设我们使用一维的每个线程块线程数或网格大小时的情况。现在我们需要迈向下一步。

\subsection{二维线程块配置}

我们需要使用二维的大小，就是这样。你需要做的仅仅是像这样设置一个 X 值，这里的 32 和一个 Y值，这就是我们需要做的
全部。那样的话我们将有两个标识符，第一个用于确定 X 维度的大小，第二个用于确定 Y维度的大小。因为这里两者都等于 32
，所以 X 维度的 \texttt{blockDim} 等于 32，Y 方向的\texttt{blockDim} 是 32。

但我想在这里提一下，这两个值的乘积绝不能超过这个值。因为我说过最大线程块大小，无论是一维还是二维都不应超过这个值，绝不能
超过这个值，即1024。如果你在你的应用中尝试使用，例如 1024 和 2，如果你将它们相乘会得到类似 2048的结果，
那么应用将无法运行。即使编译是正确的，你需要留意这一点好吗？

但只要你不使用大于这个值的数字，你可以自由调整这两个值，你可以调整它们，并且你需要这样做，因为这里的每个选项都可能影响你
应用的性能。例如你可以说64 和 16，你可以使用 1 和 512，可以使用 2 和 512，或者在 X 维度用 512，Y 维度用2。你需要测试这一点好吗？

然后对于网格大小你也可以进行选择。抱歉，这是逗号。我应该在这里加个逗号好吗？但无论如何，这很简单，你也可以在 X维度使用
一定数量的线程块，在 Y 维度使用另一个数量，并且与线程块大小类似。我们这里有两个标识符来确定 X 维度和 Y维度的大小。

例如在这个例子中，在这个实例中，我们在 X 维度有每个网格 1024 个线程块，在 Y 维度有每个网格 1024个线程块
。如果你想计算这个应用中的线程块总数，你可以将这两个值相乘。如果你将它们相乘，会得到超过 100万。每个网格的线程块数就
是这两个值，明白吗？如果你想计算这个应用中的总线程数，你需要将这两个数字相乘，其结果再乘以那个乘法的结果。将是32 乘以
32 乘以 1024 乘以 1024。这就是如果我们使用这里的配置，每个应用的总线程数。

以上只是对我们将在应用中所做更改的一个介绍。但我想从矩阵加法的定义开始讲起。简单的矩阵加法。因为我们知道每个矩阵都是二维
的，我们需要从头开始理解。因为我想告诉你们，每当我们解释二维线程和线程块的索引时，这是最困难的部分。有矩阵本身的两个维度
。

\subsection{行主序内存存储}

这并不容易，因为我们这里涉及不同的概念，有内存索引或每个元素将存储的内存地址。我们还有线程 ID 和线程块ID，但这里的
问题是我们有 X 维度的线程 ID 和 Y 维度的线程块ID。你需要理解所有这些概念是如何映射到一起的。让我们从这里开始
，我们有一个名为 A 的矩阵，需要将其与另一个名为 B的矩阵相加。结果将存储在另一个名为 C 的矩阵中，并将结果存储在
C的第一个元素中。我们需要做的是按元素依次相加，意思是我们要取第一个值，把它加到这里的第一个值，A 的第一个值加到 B的
第一个值上。这就是我们在这个例子中做矩阵加法的方式。

这只是我们需要很好理解的图表，因为我们知道这里有许多元素。但我们需要知道这些元素将如何存储在内存中，对吗？让我从这里开始，
仅以矩阵A 为例，我将只关注其中一个矩阵，同样的道理也适用于另外两个矩阵，这是我们的第一个矩阵。正如我们所见，我假设在
前两行有一些值。它们从较小值递增到较大值，我故意这样做，以确保我们是按顺序读取它们的。这只是一个假设。然后我想知道这些值
将如何存储在内存中。因为内存不是二维的，我们有一个可以在代码中通过变量定义的矩阵，我们需要知道这些值将如何存储在内存中。

内存是一维的，对吗？这个二维矩阵如何存储在一维内存中呢？它看起来是这样的，这里我聚焦于第一行，我在展示第一行将如何存储在
内存中。你需要关注的是这些值，如你所见，第一个值会在这里，下一个值在这里。它们不是以一行，而是像存储在一列中那样。但这没
关系，你需要知道的是第一行的所有值将在内存中连续存储。连续意味着每个值都紧挨着前一个值，正如你在这里看到的4、7、9、0、2、5。

第二行，第二行会以不同的方式存储，还是相同的方式？它将以相同的方式存储，这就是我们要做的。第二行将会是这样，我们不是拥有
两行，而是将所有值连续存储。就像我们在这里看到的，这是假设我们使用的是行主序（Row-major）存储。行主序意味着每一
行将在内存中连续地存储其值。我只是解释这两个术语，但我们将只关注行主序。但如果是列主序（Column-major），那么
我们将取第一列的值，它们将在内存中连续存储。矩阵A 中的二维值将像这里所示那样在内存中连续存储。

我知道我解释得比较慢，因为几分钟后事情会变得更复杂一点。总之请注意我在这里所说的，因为这非常重要。然后这里我展示了内存中
每个位置的地址。我展示这个是因为我们将在代码中需要它。我们这里有10 个位置，是的，从 0 到 9，10 个位置。我这里
假设这些内存位置的地址将从零开始，第一个位置是 0，第二个位置是1，以此类推。在地址之后，我们这里有 X 维度的值，也有
Y 维度的值，对吗？在这一边，每个元素将由一个 X 值和一个Y 值定义。但在内存中，这个值或同一个元素将仅由一个代表地址
的值定义。我们现在需要知道的是定义这些元素中任何一个的两个值X 和 Y 如何被转换或转移到这仅有的一个值，对吗？

\subsection{二维索引映射}

这就是我要解释的，我们可以说 i 或 x，这没关系。我假设我们在谈论这个值 25，这里的 25 的 X 值等于3。比如说
i 等于 3，这里的 j 等于 1。所以我们有 i 等于 3，j 等于 1。但在内存中——抱歉，在 CUDA代码中我们不
能写 3 或 1，明白吗？因为 3 或 1 并不能定义这个元素在内存中的位置，对吗？我们不能使用 3 或1。我们需要做的
是，如果我们看 25 这个值，可以看到它有一个索引8。相反我们需要地址，或者更具体更准确地说，我们需要这个元素在内存中的
索引，对吗？

内存索引是 8，但是 i 和 j 是 3 和 1，它是这样的，我们如何将 3 和 1 映射到值 8呢？这是我们使用的方程。
我们说 j 乘以行大小加上 i。把我们的值代入这里，j 等于 1，二维矩阵的行大小是5。因为我们有从 0 到 4，每行
有五个元素组成，我们将用 1 乘以 5 再加上 3，这等于8。这就是我们需要的，这就是我们将二维矩阵元素映射到内存中一个
索引值的方法。

这容易吗？这是第一步。但这里我们还没有讨论线程 ID 和线程块ID。别忘了我们有二维线程和二维线程块。我们如何将任何元素
的地址或内存索引映射到任何元素的 X 和 Y值呢？将矩阵中的元素映射到线程 ID 和线程块 ID，你看出这里的困惑了吗？
你看出这里的问题了吗？这三者之间的映射关系。

这是我们的标题：如何从线程 ID 和线程块 ID映射到内存索引。让我们先总结一下之前在章节和例子中讨论的内容。正如我们所
见，假设我们只有一个向量，就像之前讨论的那样，我们只有一个向量。这是一个一维向量，一维应用。在我们的应用中，我们假设使用
了四个线程块，这就是为什么所有的 X 线程块ID 和线程 ID 都只存在于 X 维度。

我假设这里的向量大小是这么多，我可以说是 32个元素。你可以数一下所有这些值，看看这个向量里有多少个元素。第一个元素的索
引是 0，最后一个元素的索引是31。如果我们需要这些元素的索引，就是这样。这些就是索引。我们的应用中有四个线程块，如果我
们想说 A[i] 或A[i]+B[i] 等等。

每个线程块将处理八个元素，这里的线程块 idx 是 0，这里的线程块 idx 是1。然后每个线程块由八个线程组成。正如我
们在这里看到的，每个线程块由八个线程组成。选择每个线程块八个线程并不是一个好的做法。你应该使用的最小值是每个线程块32
个线程，这等于一个 Warp 数。但这里这只是一个简单的例子。

无论如何，你还记得这个方程吗？我们曾在代码中使用这个方程来获取这些元素中任何一个的索引。它是这样的：这是线程ID。例如如
果我们看这里的这个值 3，这个值的线程 idx，在这个线程块中，线程 ID 从 0 到 7 对吗？线程 ID是 3 加上
线程块 ID。顺便说一句，这个方程和这个方程非常相似，它们非常相似。但不管怎样，让我们聚焦在这里，看到了吗？

线程块 idx，这是线程块 ID 1，我们将其乘以大小。X 维度的 \texttt{blockDim} 指的是线程块在X
维度的大小。我们知道每个线程块由八个元素组成，我们将乘以 8，那么这里的值结果将会是 11。或者 11——抱歉，不是10，是 11。
0, 1, 2, 3, 4, 5, 6, 7, 8, 9, 10, 11。我们现在使用线程 ID和线程块 ID 计算出了这个元素在这里的 ix或索引值，这就是代码看起来的样子。如果你还记得我没有在这里添加任何新代码，这只是对所有先前讨论的复习。
然后我们可以使用这里的ix 来遍历从 0 到 31 的所有元素，这里的 ix 等于 \texttt{threadIdx.x} 加上\texttt{blockDim.x} 乘以 \texttt{blockIdx.x}，
就是这样。

现在我们需要前进到更复杂的部分，即使用二维数组而不是一个向量，对吗？为了让你更容易理解，想象一下我们将使用多个向量，而不
是一个向量。这里的第一个向量，然后这里的另一个向量以及这里的另一个向量，就是这样。第一个向量如你所见，第二个向量以及第三
个向量，我希望你能看到鼠标。我不知道这里是否有什么东西能让它更醒目。

我不确定为什么这里没有可以高亮显示的东西。总之如果我们有一个向量，我们只使用这个方程来识别该向量内的任何元素。但如果我们
有多个向量，我们可以使用另一个方程，这是针对X 维度的。类似的还有一个针对 Y 维度的，对吗？如你所见，ix 将完全是我
们这里使用的那个，但这仅适用于第一个向量。

如果我们有多个向量，我们可以使用 iy 来定义我们正在使用哪个向量，对吗？这是用来确定向量内的任意元素，而 iy是用来确
定我们正在使用哪个向量。这就是我们可以从向量过渡到矩阵，或者说从一维过渡到二维的方法。这是用来确定一行中的哪个元素，而这
是用来确定我们具体在处理哪一行。这里如你所见，我使用ix 来确定向量内或具体到某行内的任意元素。因为如果是二维矩阵，那么
我们需要行和列，这完全是一个与这个方程类似的方程，对吧？所以这是一个类似的方程，用于确定二维数组中的哪个元素。

ix 是我们从这里计算的，iy是我们从这里计算的。这两个值定义了矩阵中这里的任何元素。例如这里的这个元素由我们从此处计算
的 ix 值和从此处计算的iy 值定义。我知道这仍然很复杂，因为我们现在只讨论了线程和线程块ID。但是内存索引或元素索引
呢？这是元素索引，我们最终可以使用的元素索引好吗？就像这里我们将在这里使用它，对吗？代替 X值，这里我们是从线程 ID、
线程块 ID 映射到元素索引，明白吗？但我想问的是矩阵的大小具体体现在哪里？和线程块ID、元素索引。因为我们涉及三个部分：
线程，这个是——抱歉。

最后一部分是矩阵本身的大小，让我带你看一下这个。正如我所说，我们需要理解三个部分，而这只是一个实际例子，用来理解和想象这
些部分是如何映射到一起的。在这个例子中，矩阵大小是从0 行到 5 行。这个矩阵有六行，这是矩阵 A，我们在讨论这个矩阵的
大小，我们有六行和 8 列，因为列从 0 到7。是八列，记住这个好吗？八列和六行。

\subsection{二维矩阵配置实例}

在这个例子中我们有多少个线程块？线程块的数量你可以在你的应用中配置。例如这里我们有——让我在这里展示矩阵大小 X 和 Y
值，我们有六行和 8 列。

我在这里假设我有六个线程块，并且它们是二维大小的。对于二维线程块和二维线程索引，我们有二维网格大小。网格大小指的是每个网
格的线程块数，我们讨论的是线程块的数量，在X 维度我们有三个线程块对吗？一——抱歉，这个应该在 Y 维度，不是三，不是
X 维度好吗？不管怎样，所以我们有1、2、3 好吗？我们有三行线程块，而在垂直维度，总之我们有两个。这是一，这是 2。

这个例子中的线程块总数是六个，对吗？这就是它们如何映射到这个矩阵中的不同元素。如你所见，线程块 (0,0)将处理这里的这
些元素，这八个元素好吗？你现在不需要看下一步就可以明白，每个线程块将由八个线程组成，因为它处理八个元素，对吗？这适用于所
有线程块好吗？

第二个线程块，即线程块 (1,0)——X 维度的 1 和 Y 维度的0，它也处理八个元素，这就是二维线程块如何映射到这个
矩阵中的所有元素。然后我们有二维线程块大小好吗？我们说总线程数将是2 乘以 4 等于八个线程。每个线程块的线程数是 2
乘以 4，因为这里每个线程块有两行和 4 列。

我们可以说这里的 \texttt{blockDim.x} 是 4，\texttt{blockDim.y} 是 2。

\texttt{blockDim} 指的是线程块大小。你还记得 \texttt{blockDim}是什么意思吗？如你所见
，我们在 X 维度有 4，Y 维度的线程块大小是 2。只写 \texttt{dim.x} 而不写\texttt{blockDim.x} 是错误的。
因为我们还有网格大小的 dim，我们必须在这里明确，我们有\texttt{gridDim} 和\texttt{blockDim}，像这样写是错误的。但不管怎样，\texttt{gridDim} 
在 X 维度和 Y维度将是 3 和 2，然后是内存索引。

\texttt{gridDim} 和 \texttt{blockDim} 将分别在 X 和 Y 维度是 3 和2，这里是
矩阵大小，这是二维线程块和线程，这是内存地址。我不喜欢说内存地址，但我只是为了让你们更容易理解，我在每个元素内添加了内存
索引。但我们可以说内存和这些索引，好的。例如第一个元素的索引将是0，第二个元素是 1，第一行的索引将从 0 到 7。这就
是它们在内存中的排列顺序。

如果我们看第一个元素它是 0，最后一个元素它是 47，我们可以理解这个矩阵有 48 个元素，对吗？然后继续索引8、9、10，以此类推，直到我们到下一行 16、17 等等。元素总数是 48 个，它们从 0 到 47索引。但我们没有每个元素的具
体值，好吗？我想让你们知道我们如何在这三个部分之间进行映射。我们不需要关注值本身，让我以这个例子为例，这个值第44 个元
素，我将遮住其他线程块，以免用不同的东西让你们困惑。

我将以这一个为例，这是我们在处理的线程块，这些是你们正在处理的元素，这是 ix 值。记住我们需要计算ix、iy，然后是索
引，我们已经知道索引是 44，但在这里我是为了验证原理。我是为了证明是的，它是44，并且这个值可以使用这里的这些方程计算
出来，这是 ix 值。

我想问你，你能从这个线程块中给出 \texttt{threadIdx.x}吗？我们知道每个线程块由八个线程组成。这些线程
有两个维度，一个在 X 一个在 Y。你能给出这里\texttt{threadIdx.x} 的值吗？我们只讨论 idx，不是两者好吗？你能给我这个元素的\texttt{threadIdx.x} 吗？是的，请仔细想一想好吗？

我在这里直接给你这个值，如你所见它是 0，为什么？因为我们从这里开始，这是 0、1、2、3，这是针对 X 维度。对于 Y
维度这一行将是 0，这一行将是 1，对吗？如果我们看 X 维度这是 0、1 和2，但这是针对两个维度的，不是针对一个维度。

现在你需要考虑 \texttt{blockIdx.x}，记住这是线程块 ID，我只需要 idx 即 X 维度的ID。如果
你看这个值，你会发现它是 1，因为这是 X 维度上的值。我们可以从这里得到这个值好吗？然后我们需要知道 X维度
的 \texttt{blockDim}，即 X 维度的线程块大小 dim。X 维度的\texttt{blockDim}
是 4 对吧？这里的值 ix 值将等于 4 乘以 1 加 0，所以等于 4。

现在我们有了 ix 值，需要计算 iy 值对吗？是的，是 1 对吧？你能告诉我同一个元素的线程 ID 吗？但是在 Y维度
不是 X 维度。是 0 和 1，这个值是 1。你能告诉我 Y 维度的线程块 ID 吗？不是 X 维度。然后你能告诉我Y
维度的 \texttt{blockDim}，即 Y 维度的线程块大小吗？它会是——线程块 ID 像是 2 对吧？

1、2。从这个方程得到的总值将是 5，这是 ix 和iy。现在我们需要计算索引本身，即我们将在代码中使用的索引，对吗？A
的索引加上 B 的索引等于这个值将存储在 C的索引中，对吗？那么让我在这里展示索引值，这里的 iy 是 5，所以这是 iy 乘以 nx。你能告诉我 nx 吗？


这完全来自这里矩阵大小明白吗？因为这是 nx 即矩阵 A 在 X维度的大小，如果是这样，我们可以得到值的列的数量。因为每
行——这里每行如果你还记得每行有八个元素或八个值，每行的大小是8，这就是 nx 的值。你现在明白了吗？线程 ID 被转换
为 idx值，而这里我们使用的是矩阵大小，这就是这里三个部分如何映射在一起。总之我们可以说这个矩阵在 X 维度的大小是8
，这正是我们一开始就在这里看到的那个值。5 乘以 8 加上 ix 值即 4，总值将是 44。

我希望这能让你更容易理解和想象二维线程是如何映射到我们正在处理的矩阵的任何元素的。这只是所有这些内容的一个总结。我们需要
这些方程的结果，结合以列数表示的矩阵大小来得到内存索引。就像我们在这里看到的。我们需要线程 ID 和线程块 ID 来获得
ix 和 iy。线程 ID 和线程块 ID 来计算 ix 和 iy，这只是两句话来总结所有这些。

这是 CUDA 编程最优秀的书籍之一。这些例子或这个映射，我想对你们说实话，并说明这是从这本书中引用的，而这正是我希望你
们看的，它会带你一步一步学习。我希望你们看这本书并继续学习。其他例子我想是在第三章，第三章或第二章我记不清了，但我就是从
这本书中引用的。这与我们将看到的命中和未命中数量有关，但我现在不打算解释这个。让我带你们看我们的代码。在去看我们的代码之
前，我想在这里给你们看些别的东西，让我给你们看些别的东西。

如果你在处理二维线程和线程块，你可以定义 \texttt{blockDim}。就像你在这里看到的\texttt{blockDim}，这只是一个变量名好吗？然后你可以使用这里的方程来计算每个维度中剩余的线程块数量。但你可以将其定义
为32 乘以 32。如我们讨论的 X 维度每个线程块 32 个线程，Y 维度每个线程块 32个线程。因为我想告诉你，每
个矩阵都有特定数量的元素，对吧？例如假设我们有一个矩阵，大小是 X 维度 4096 个元素乘以Y 维度 4096 个元素
，那么元素总数将是这两个值相乘的结果。

元素总数必须等于你将在应用中使用的总线程数。我们这里有两个配置：每个线程块的线程数和线程块数量。如果你像我们在这里做的那
样，手动确定每个维度的线程数，比如32 和 32，那么你必须使用一个方程来计算线程块的数量。你不能假设这里是 32 乘以 32，
而那里例如是 1024乘以 1024。如果你只使用一个维度，你可以使用这里的这个方程，这就足够了。不，这不是它的工
作原理，因为你需要确保所有维度中所有值的乘积将给出一个线程数，等于我们相加的矩阵中的元素数量。如果你使用两个维度，你可以
在这里使用这个逗号来为下一个维度或Y 维度设置值。如果你运行这个应用，并且如果你在这里使用相同的矩阵大小，当你尝试使用例
如 NCU分析时，你会在输出中看到这个。我将向你展示这个例子，我将运行它，但我只是在展示代码前解释这个概念，因为代码不是
困难的部分，理解代码背后是什么，这个代码背后的原理是什么，才是你需要知道的最重要的事情。当你尝试分析这个具有这些配置的应
用时。

以及这个矩阵大小。你会在输出中看到这个，你会看到我们的 \texttt{blockDim} 是这里的这个指令给出的 32 和 32。并且从这个方程你会看到，我们在 X 维度有 128 个线程块，在 Y 维度有 128个线程块。我想问你为什么是
128？为什么不是例如 256？

为了解释这个，我想告诉你，我们必须看一下这里的这个图表。我告诉过你，从这个输出中我们有 128 个线程块。正如我们从0、1、2 
直到 127 所看到的，这里的每一行代表线程块的数量。列也是如此，从 0、1、2 开始直到127。然后我想问你，
你能告诉我 X维度有多少线程吗？这个图表代表线程块的数量，并且对于每个线程块，正如我们在这里看到的，每个线程块的大小是
32 乘以32，对吧？我们有 128 个线程块，每个线程块在 X 维度有 32个线程。这很简单，你只需将这个值乘以这个值，
你就能得到 X 维度的这个值；如果你将这个值乘以这个值，你就能得到这个值 Y维度。这就是我要告诉你的，你需要确保 X
维度和 Y 维度的总线程数应该等于矩阵大小。

我不打算在这个文章中讨论这个。我将在下一个文章中讨论它们。我希望你看一下代码，它是这样的，我在这里要做的是复制这个，然后
转到我们的新例子。

\subsection{一维与二维代码对比}

但在看今天讨论的代码之前，我想带你回顾一下第一个维度——抱歉，一维线程和线程块。这里我们称它为\texttt{vectorAdd} 并找出差异。我需要你在这里看一下内核，GPU 内核它由
\texttt{\_\_global\_\_ matrixAdd} 开始。如果你看这个我们有 ix 和iy，你可以给它们起任何名字，
你可以叫 row、col 或者你希望的ix、iy。但这里我们只有一个，因为这是向量，向量不是矩阵。这里我们说 C[ix]
等于 A[ix] 加上B[ix]。这里我们必须使用行值和列值重新计算索引，这是唯一的区别，没有其他区别。

好的，然后在代码内部，如你所见，我们通过将 nSize 和 mSize 乘以 float的大小来确定每个矩阵的大小。因为
你在处理单精度值，这里的大小是通过两者相乘来定义的，而不是像这里那样大小只乘以一个维度。N和 M 在这里定义，N 是 4096，M 是4096。让我看看这里，看这个只有一个维度的大小，因为这是向量，而这是矩阵，对吧？一旦你定义了每个矩阵的总
大小，你将在CPU 端为它们分配空间。

然后初始化 A 和 B 并在 GPU 端在 GPU 内存中分配一个相似的空间。然后将 A 和 B的数据从主机复制到设备，
或者说从 CPU 复制到 GPU。正如你在这里看到的，从主机到设备，接着我们将定义 X 和 Y维度的线程块大小，以及 X 和 Y 维度的网格大小。
正如你在这里看到的，我只是想让你看看 X维度的这一部分，你会发现它是相等的，它非常相似，和这里的
这个方程这完全是同一个方程。但这里我们只使用一个，没有逗号，因为这是向量不是矩阵，所以我们只需要一个维度。

顺便说一下，在这个例子中，你也可以使用一维的线程块大小和网格大小来进行矩阵加法。这就是为什么我告诉你看这本书，因为它有这
个例子，我无法在这里解释所有例子，但我需要给你一些任务。它将会工作好吗？与其使用两个维度来进行两个矩阵加法，尝试使用一个
维度来相加这两个矩阵。并且你可以研究这两种情况下的性能。总之这里这个我只是想跳过，因为它涉及 L1缓存，所以今天我不想讨
论它。你不需要添加它好吗？你可以移除它。

然后我们在这里启动内核来执行。我们发送 A、B、C 以及 X 和 Y 维度的大小。一旦这里的全局执行完成，它将在这里打印
A、B 的前两个元素及其值。

以及 C 以检查结果是否成功计算出来，这就是整个代码。我们需要做的就是来到这里，这个文件叫做\texttt{two\_dimensions\_01.cu}。
我在这里要做的是来到这里使用 NVCC编译这个文件。我们将把输出或可
执行文件保存为这个名字，然后点击回车，这里没有错误，我们需要做的就是点击\texttt{./}，然后是 \texttt{.exe}，
然后点击回车，正如我们在这里看到的这等于 16 点多。

第二个值同理，因为我们在这里打印这两个值来验证我们工作的正确性，这是输出 8.4 加7.8，这就是它的样子。但同样你也可
以来到这里。我需要你看一下这一部分，这是我们的命令，并且写例如\texttt{ncu} 来查看性能分析的输出。正如我们看
到的。你还记得幻灯片上的这个吗？这里这个来自幻灯片。

计算出的两个维度的线程块数量是 128 和 128 个线程块。如你所见，线程块大小的配置是 X 维度 32，Y 维度32。
然后你可以查看 SM 的利用率，可以查看不同内存级别的吞吐量，如 DRAM、L1、L2 等等。你会发现线程块大小是1024 对吧？同样你也可以查看启动统计信息，线程块大小是 1024，为什么？
如果你将这些值相乘，32 乘以32，这将是 1024。因为每个线程块的总线程数将是 32 乘以 32，然后网格大小是 16384。

如果你将这个值乘以这个值，这像是 1600万，对吧？这是总线程数。因为我们有线程块大小和网格大小，请记住这一点，我将来这
里计算 4096 乘以相同的数字。这将是1600 万，与总线程数相同，元素总数。因为这是矩阵大小，对吧？矩阵大小中的所有
元素，它们的总数应该等于我们应用中的总线程数。

\subsection{性能分析与优化}

占用率我们有 50\% 的占用率，受线程块限制的线程数和这里的线程块。然后每个 SM的波次（Wave），我已经在前面的章
节中解释了波次是什么意思，所以我们需要再次复习。我们每个 SM同时只能运行一个线程块，这非常糟糕。原因是我们在每个线程块
中使用了太多的线程数。你使用了每个线程块的最大线程数，这等于32 个 Warp。但是如果你将这个值从 32 乘以 32
改为更小的值，你会看到这个值应该会增加。这里的每个线程块由32 个 Warp 组成，这就是为什么我们不能使用超过一个线程
块。因为每个线程块都很大，每个线程块有 32 个Warp。让我试试这个，例如让我们来到这里。写成 16 乘以 16，然后
编译并分析性能。

受线程块限制的线程数数量，这里是6。如你所见看这个，这意味着我们可以使用——我们取这里的最小值，对吧？我们可以同时运行每
个 SM六个线程块，将这个值与之前的值比较好吗？我不知道这是否会影响总执行时间，但我们可以检查一下，让我们比较一下执行时
间。这里的占用率是74\%。

因为我们有多个线程块，每个都有多个 Warp，这是 74\%。但在第一种情况下，你可以看到只有50\%，这些是性能指标，
而这是分析这些指标的方法。我这里把执行时间存储在这个文件中，我在前面的一个章节中解释过这个。

这将在这里给我们一些性能指标，我在这里需要做的就是写 \texttt{sh}，然后点击回车，这里的执行时间大约是 478
微秒。例如周期数、执行时间，我们只需要看这个以及 L1 和 L2 的命中率，L1 命中率大约是 1\%。几乎是0\%，L2 
也几乎是 0\%，但这里另一个 L2 命中率等于33\%。我将在下一个文章中解释它们，这是我在本课程中最喜欢的文章之一。

下一个。总之我们需要关注这两个值：时间和活动周期。然后回到这里的代码，将这些值恢复为 32 乘以32。保存、编译。你必须
编译好吗？然后再次运行相同的脚本。

让我们看看这里。是的，如你所见，这种情况下的执行时间与这种情况相比较长，我想现在这很容易理解了。这里的占用率是 73\%
或 74\%，这里的占用率是 49\% 或 50\%。然后我认为不同级别缓存的命中率是相同的，33\%。

这就是今天的全部内容，我希望这对你来说很容易理解。这是我们进入本课程下一阶段或第二阶段的开始。今天就到这里，我们下个文章
见。


